\documentclass[letterpaper,journal]{IEEEtran}
\usepackage{amsmath,amsfonts,amssymb}
\usepackage{algorithmic}
\usepackage{algorithm}
\usepackage{graphicx}
\usepackage{cite}
\usepackage{booktabs}
\usepackage{multirow}
\usepackage{stfloats}
\usepackage{url}
\usepackage{xeCJK}
\usepackage{textcomp}
\usepackage{caption}

\xeCJKsetup{CJKecglue=}
\setCJKmainfont[BoldFont={Noto Serif CJK SC Bold}]{Noto Serif CJK SC}
\setCJKsansfont{Noto Sans CJK SC}

% IEEE 风格 Caption
\captionsetup[figure]{font=small,labelsep=period}
\captionsetup[table]{font=small,labelsep=period}

% 紧凑型表格（IEEE 推荐 \footnotesize）
\renewcommand{\arraystretch}{1.15}
\setlength{\tabcolsep}{3pt}

\hyphenation{op-tical net-works semi-conduc-tor IEEE-Xplore}

\begin{document}

\title{SPA-KLT: 一种面向光流视觉惯性里程计的自适应稀疏光流跟踪算法}
% If English title needed:
% \title{SPA-KLT: An Adaptive Sparse Optical Flow Tracking Algorithm for Optical-Flow Visual--Inertial Odometry}

\author{
  第一作者,
  第二作者,
  第三作者%
  \thanks{稿件准备中，请补充作者单位、通信作者邮箱及基金资助信息。}
}

\markboth{IEEE Robotics and Automation Letters, Vol.~XX, No.~XX, XXXX 2026}
{Author: Title}

\maketitle

\begin{abstract}
视觉惯性里程计（visual-inertial odometry, VIO）是机器人、增强现实与自主移动平台实现位姿估计的关键模块。基于 Lucas--Kanade（LK）稀疏光流的间接法因轻量与可移植性在嵌入式平台获得广泛部署，但其固定搜索窗与金字塔层数导致算力浪费、IMU 高动态下特征位移远小于搜索窗、高速旋转与运动模糊下前端鲁棒性退化、回环校验缺少前置几何筛选，长期构成难以根治的四大结构性短板。半直接法中广泛使用的稀疏光度对齐（sparse photometric alignment, SPA）虽能提供全局运动先验，却依赖小块 patch 重叠、单位深度投影、亮度一致与强响应纹理等强假设的联合成立——仅在帧间运动适中、纹理充足、曝光受控、视场受控的工况下近似可用，其他工况下假设失效会反向拖累前端。

\textbf{本文的核心创新是提出 SPA-KLT}（\textit{Sparse Photometric Alignment for Adaptive KLT}）：一种面向光流 VIO 系统的\textbf{条件性可用增强层}。其将 IMU 预积分旋转作为稀疏光度对齐的迭代初值，并以对齐残差 $\chi^2$ 作为下游 KLT 搜索窗与金字塔层数的自适应依据，对齐失败时自动回退至固定参数 KLT，从而把``粗定位—精跟踪''两阶段耦合在一起；并在此基础上引入仅在稀疏对齐成功路径上启用的稀疏独占节流策略（关键点检测间隔触发与 RANSAC 置信度自适应），以及基于稀疏旋转的回环几何门机制。

\textbf{本文的结论}：以``适用场景''作为前置判据——适用场景（帧间运动适中、纹理充足、镜头畸变受控、曝光稳定、强响应特征点分布）下，前端耗时平均降低 $20\%$ 以上、对齐模块净收支比达 $2.13$、自适应最低档单帧像素操作量降至 baseline 的约 $3\%$，KLT 不再是前端主要瓶颈，长程累积漂移显著抑制；不适用场景（大像素位移、低纹理走廊、宽视场、自动曝光剧烈跳变、相机剧烈抖动等）下 SPA 强假设失效，方法以``$\chi^2$ 阈值升档与对齐失败即回退''形式保证不反向拖累前端；该局限与稀疏对齐方法族（含 SVO 等）共享，需架构级改变才能进一步突破。

\textbf{本文解决的工程问题}：以零侵入方式集成至 VINS-Mono，作为通用增强层不改动后端 Ceres BA、IMU 预积分与位姿图优化代码，为光流 VIO 前端提供一条可量化的``条件性可用''工程化路径。
\end{abstract}

\begin{IEEEkeywords}
视觉惯性里程计；稀疏光度对齐；自适应 KLT 光流；前端跟踪鲁棒性；EuRoC 数据集；旋转几何回环门
\end{IEEEkeywords}

\section{引言}
\IEEEPARstart{视}{觉}惯性里程计（visual-inertial odometry, VIO）是机器人、增强现实（augmented reality, AR）与自主移动平台实现位姿估计与建图的关键技术之一~\cite{cadena2010past,davison2007monoslam,klein2007ptam,engel2014lsd,mur2015orb,forster2014svo,qin2018vins}。按照对图像信息的利用方式，现有 VIO 方法大体可分为两类技术路线：其一为基于特征的间接法（如 ORB-SLAM~\cite{mur2015orb,campos2021orb}），先提取并匹配稀疏特征点，再以重投影误差作为视觉残差项送入后端优化；其二为基于像素灰度的直接法（或半直接法），直接利用图像像素灰度构建光度残差估计帧间运动~\cite{engel2017dso,forster2017svo}。在间接法体系内部，Lucas--Kanade（LK）稀疏光流~\cite{baker2004lucas}作为 VINS-Mono~\cite{qin2018vins}等代表性系统的前端支柱，以``小运算量、强可移植性''在嵌入式平台获得了广泛部署。

然而，LK 光流路线的固有局限同样不容忽视：固定搜索窗与固定金字塔层数导致算力浪费、IMU 高动态时特征位移远小于搜索范围、对高速旋转与运动模糊场景缺乏鲁棒性。另一方面，回环检测作为消除长距离轨迹累积漂移的核心环节，长期依赖词袋（bag-of-words, BoW）模型对图像外观进行检索~\cite{galvez2012bags}，在重复纹理场景与视角变化下容易引入虚假回环，而现有方案缺少基于运动一致性的前置筛选。

具体而言，基于光流的 VIO 体系在实践中存在以下四方面固有局限：

\begin{enumerate}
\item \textbf{固定参数 KLT 光流缺少全局运动先验。} 经典间接法前端普遍采用固定搜索窗与固定金字塔层数的 LK 光流跟踪，在弱纹理或快速运动场景中极易丢失特征；当 IMU 处于高动态时，帧间特征位移可能远小于搜索窗范围，大量计算被浪费在外点判别上；与此同时 KLT 本身又是最耗时的前端环节，算力浪费与鲁棒性不足并存。
\item \textbf{前端跟踪鲁棒性在高速旋转与运动模糊场景下退化。} 当载体经历快速旋转或图像出现明显模糊时，固定参数 LK 极易失锁；特征丢失后又被默认检测节奏填补，进一步占用本就紧张的算力预算，形成``算力—鲁棒性''的恶性循环。
\item \textbf{回环校验缺少前置筛选环节。} 现有方案大多仅在词袋检索后直接执行重投影误差校验，缺少基于运动一致性或几何一致性的前置门限，重复纹理场景下的虚假回环难以被抑制，从而引入错误的全局约束，造成长程漂移难以收敛。
\item \textbf{稀疏光度对齐的可用场景具有显式边界条件。} 半直接法中广泛使用的稀疏图像对齐（sparse image alignment）以帧间像素灰度一致性为代价函数，其有效性依赖若干强假设——帧间位移需落在 ``小块 patch 重叠范围'' 之内、相机近似满足单位深度投影假设、图像亮度近似恒定或可被低阶仿射项吸收、特征点邻域内纹理可被光度一致性约束可靠区分。这四项假设同时成立的近似场景是 ``持续飞行 + 帧间运动足够 + 充足纹理 + 受控曝光'' 的 MAV 类数据集；其他场景下，假设的系统性偏离会使稀疏对齐的精度退化甚至失效——这一点在 SVO 原文已明确指出，原作者在多数据集评测中报告稀疏对齐在低动态、宽视场、室内重复纹理、自动曝光、相机剧烈抖动等条件下相对直接法或 VIO 系统的提升幅度有限甚至反向~\cite{forster2017svo}。本文在 VINS-Mono 框架内将稀疏对齐定位为``条件性可用增强层''而非无条件可靠的视觉前端，并据此设计适用场景识别、降级容错与稀疏独占节流门控机制，避免在不适配场景中反向拖累前端。
\end{enumerate}

针对上述局限，本文提出 \textbf{SPA-KLT}（\textit{Sparse Photometric Alignment for Adaptive KLT}），一个面向光流 VIO 系统的通用增强层。其核心思想是将 IMU 预积分得到的帧间旋转作为稀疏光度对齐迭代初值，以对齐残差 $\chi^2$ 作为下游 KLT 参数选择依据，从而把``粗定位—精跟踪''两阶段耦合在一起。本文主要贡献如下：

\begin{enumerate}
\item \textbf{基于 IMU 旋转约束的稀疏光度对齐自适应 KLT 跟踪模块。} 将 IMU 预积分得到的帧间旋转作为光度迭代初值，构造 IMU 正则化的光度残差函数，并将对齐残差 $\chi^2$ 作为 KLT 搜索窗（$5{\times}5$ / $9{\times}9$ / $15{\times}15$ / $21{\times}21$）与金字塔层级（1/1/2/3 层）的自适应选择依据；设计降级容错机制，对齐失败时自动切换至固定参数光流。
\item \textbf{两项稀疏独占前端节流策略。} 在对齐成功的前提下启用关键点检测间隔触发（每 3 帧检测一次）与 RANSAC 置信度自适应（$0.99 \to 0.95$），前者将特征检测开销由 $0.60\,\mathrm{ms}$ 降至 $0.20\,\mathrm{ms}$/帧，后者将离群抑制开销降低约 $40\%$，三项合计实现对齐模块净收支比 $2.13$。
\item \textbf{回环候选帧稀疏旋转几何门机制。} 复用前端稀疏相对旋转与 DBoW2 候选位姿所蕴含的相对旋转进行一致性比对，超出阈值则在前置阶段直接拒绝候选，为复杂场景下的鲁棒性提供第二道保险。
\item \textbf{在 VINS-Mono 上零侵入集成并完成 EuRoC MH\_01--MH\_05 与 V101--V103、V201--V203 共十一个序列的消融与对比实验。} SPA-KLT 在十序列上优于或持平基线（ATE 平均 0.95x），前端总耗时平均降低 $21.4\%$（V201\_easy 序列改善最大达 $45.8\%$），MH\_04 困难序列 RMSE 相对基线降低 $13.2\%$；在 V203 difficult 等极端场景下方法存在固有局限，末节专项分析失效根因与改进方向。
\item \textbf{面向极端困难场景的系统性局限分析。} 通过 V203 difficult 序列的专项实验，揭示稀疏光度对齐在大像素位移、快速旋转与低纹理场景下的失效机理（逆合成公式假设失效、IMU 先验偏离），分析 ON/OFF 从 $2.33$→$1.70$ 但仍达 $1.70$ 的结构性差距，为后续架构级改进提供方向性指引。
\end{enumerate}

\textit{与本文方法相关但定位互补的近期代表性工作。} Liu \emph{et al.} 的 SD-VIS~\cite{liu2020sdvis} 在 EuRoC 上以 ``keyframes 走 KLT、非-keyframes 走直接法'' 的二分流水线实现单帧跟踪 $6.72\,\mathrm{ms}$；Lin \emph{et al.} 的 LET-NET~\cite{lin2024letnet} 针对 LK 光流亮度恒定假设在 HDR 下失效提出了轻量 CNN 增强并在 VINS-Mono 上集成报告 translation error 降低 $93\%$；Chen \emph{et al.} 的 OTPL-VIO~\cite{chen2026otpl} 以深度线段描述子与熵正则 optimal transport 缓解低纹理与光照剧变下点特征不稳定问题；Qiu \emph{et al.} 的 DSInit~\cite{qiu2026dsinit} 以 prediction function 把 stereo VIO 初始化从 ``特征追踪 $+$ SfM'' 多阶段链改造为直接光度优化，初始化总耗时降至 $12.42\,\mathrm{ms}$。上述工作分别从 keyframe 划分、亮度不变性、特征替换与初始化直化路径切入，与本文的 ``IMU 旋转约束引导的稀疏光度对齐 $+$ KLT 自适应调度'' 互补但不重叠，详见 §II 分析。


\section{相关工作}

\subsection{视觉前端光流优化方法}
经典间接法视觉前端普遍采用固定搜索窗与固定金字塔层数的 Lucas--Kanade（LK）稀疏光流~\cite{baker2004lucas}进行特征跟踪。围绕固定参数 LK 光流的不足，已有研究提出多种改进方案，例如基于图像金字塔的自适应窗口调整、双向光流校验、逆向光流一致性约束、特征点亚像素重采样等~\cite{shi1994goodfeatures}。然而，上述方案均沿用``先提取特征、后光流跟踪''的传统流水线，未引入像素级光度对齐提供全局运动先验。

\subsection{半直接视觉惯性状态估计方法}
以 SVO~\cite{forster2014svo,forster2017svo}为代表的半直接法在前端融合稀疏特征跟踪与像素级光度对齐。SVO 的光度对齐模块仅用于局部位姿微调与深度滤波器初始化，并非作为前端光流的先验引导；此外 SVO 缺少自适应算力调度策略。基于 SVO 思路的扩展工作（如 SD-VIS~\cite{liu2020sdvis}、点线半直接 VIO~\cite{gao2022sdplvio}等）多聚焦于深度估计与回环闭合，未把光度对齐的旋转输出作为回环几何一致性校验的输入信号。在 ``嵌入式/光照变化场景的 LK 光流增强'' 这一更狭义目标下，Lin \emph{et al.} 的 LET-NET~\cite{lin2024letnet} 在传统金字塔 LK 框架中嵌入由轻量 CNN 提取的 ``光照不变特征图'' 作为光度项输入，并以 CNN-score map 替代 Shi-Tomasi 关键点检测，在 CPU 上以 $190\,\mathrm{Hz}$ 推理并集成到 VINS-Mono 后于公开 HDR 数据集报告 translation error 降低 $93\%$；该工作可视为 ``亮度恒定假设失效'' 维度上的 LK 光流学习化增强，与本文 ``KLT 自适应调度'' 维度互补但不重叠——前者替换光度项的定义，后者保留灰度假设但调整搜索参数。

\subsection{回环检测与校验优化方案}
回环检测的主流方案基于 DBoW2~\cite{galvez2012bags}等词袋模型对图像外观进行检索，并辅以时间一致性、几何一致性等多级校验~\cite{mur2015orb,campos2021orb}。多数回环优化方案仅依赖特征几何约束，无法复用前端光度匹配信息做前置筛选，对重复纹理场景下的虚假回环抑制能力有限。

\subsection{面向资源受限设备的轻量化 VIO}
面向微型飞行器、智能眼镜等嵌入式平台的轻量化 VIO 路线与本文关注的``光流 VIO 前端效率''问题紧密相关。K\"uhne \emph{et al.} 在 LEVIO~\cite{kuhne2026levio} 中以 ORB 特征 $+$ 8$-$point/EPnP 紧耦合 $+$ GAP9 多核并行化 $+$ 自适应 LM 优化器为完整管线，在 EUROC MAV 数据集上将全 6$-$DoF VIO 流水线压缩到 $<\!100\,\mathrm{mW}$、20 FPS，并显式给出与 VINS-Mono、ORB-SLAM3 在标准分辨率与 QVGA 两种工况下的 ATE 对照表。该工作在工程层面与本文关注同一目标阵地（嵌入式光流/特征 VIO），但所采取的路线是``特征点法 $+$ 硬件并行化 $+$ 自定义优化器''，与本文的``基于 LK 光流的算法自适应调度''路线正交——可视为轻量化 VIO 两条互补路径的典型代表。

\subsection{VIO 初始化与特征形式增强}
除运行阶段的视觉前端之外，近年来亦出现两个与本文方法存在动机同源性但作用阶段不同的改进方向。\textit{（i）将 VIO 初始化从多阶段链直化为像素灰度优化}。Qiu \emph{et al.} 的 DSInit~\cite{qiu2026dsinit} 在 stereo VIO 初始化阶段以 prediction function 直接建立 IMU 参数与像素坐标的映射，最小化稀疏点的光度误差，省去 ``特征追踪 $+$ 视觉 SfM $+$ 视觉惯性对齐'' 多阶段中间量；在 EuRoC 与 TUM-VI 上 3-frame 初始化成功率从 VINS/SNEC 的 $15.5\%$--$42.6\%$ 提升至 $58.8\%$--$100\%$（V203 上），10-frame 初始化总耗时降至 $12.42\,\mathrm{ms}$（基线约 $14$--$16\,\mathrm{ms}$）。该工作与本文共享 ``以原始像素灰度替代带噪中间量'' 的动机，且明确指出 ``运动模糊下亮度恒定假设违反是其直接法的已知边界''，与本文 §VI 关于 V203 失效根因的分析一致。\textit{（ii）以线特征替代/补充点特征以缓解低纹理场景}。Chen \emph{et al.} 的 OTPL-VIO~\cite{chen2026otpl} 在 stereo VIO 中引入 ``沿片段采样的深度线描述子 $+$ 熵正则化 optimal transport 全局匹配 $+$ 长度自适应可靠性加权''，在 EuRoC MH\_04/MH\_05 与 V103/V203 等 hard 序列上提升最显著（EuRoC 平均 ATE $8.06\,\mathrm{cm}$ vs AirSLAM $11.18\,\mathrm{cm}$，提升 $27.9\%$；UMA-VI 光照变化序列平均 $25.5\,\mathrm{cm}$ vs $44.1\,\mathrm{cm}$，提升 $42.2\%$），与本文关注的 ``低纹理 $+$ 光照剧变 $+$ 强响应特征稀缺'' 场景形成强互补。两条路线与本文的关系可概括为：DSInit 在系统起始阶段替代 ``特征追踪 $+$ SfM'' 中间量，本文在稳态跟踪阶段替代 ``固定参数 LK''；OTPL-VIO 在特征本体维度补充 ``点特征不足时的几何约束缺失''，本文在跟踪参数维度补充 ``固定参数 LK 的算力浪费''——分别从初始化、跟踪、特征本体三个独立维度对稀疏光流 VIO 前端进行增强。

\section{基础理论铺垫}
本节介绍 SPA-KLT 方法所依赖的通用理论工具。符号约定：局部坐标系（IMU 体坐标系 $(\cdot)^b$）与世界坐标系 $(\cdot)^w$，相机坐标系 $(\cdot)^c$ 与 IMU 坐标系的外参 $(p_{cb}, q_{cb})$ 假定已离线标定完成。

\subsection{IMU 预积分模型}
IMU 输出的加速度计与陀螺仪原始测量在体坐标系下表示为
\begin{equation}
\hat{a}_t = a_t + b_{a,t} + R_{t}^{w} g^{w} + n_{a},\quad
\hat{\omega}_t = \omega_t + b_{g,t} + n_{\omega},
\label{eq:imu_meas}
\end{equation}
其中 $b_{a,t}$、$b_{g,t}$ 为零偏并建模为随机游走，$n_{a}$、$n_{\omega}$ 为高斯白噪声。在相邻两帧图像 $b_k$ 与 $b_{k+1}$ 之间对陀螺仪测量进行预积分，可得到帧间旋转增量
\begin{equation}
R_{\mathrm{imu}} = \prod_{i} \exp\!\left((\hat{\omega}_{i} - b_{g,i})\, \Delta t_i \right),
\label{eq:imu_preint}
\end{equation}
其中 $\Delta t_i$ 为第 $i$ 段 IMU 积分步长。该预积分旋转将作为稀疏光度对齐的初值先验使用。

\subsection{金字塔 Lucas--Kanade 光流}
稀疏光流跟踪以特征点邻域像素块的灰度一致性为基本假设。给定上一帧中第 $j$ 个特征点的像素坐标 $x_{j}^{(k)}$，求解当前帧像素坐标 $x_{j}^{(k+1)}$ 的标准 LK 迭代形式为
\begin{equation}
\delta p^{(n)} = \left( \sum_{x\in W} G^{\top} G \right)^{-1} \!\! \sum_{x\in W} G^{\top}\!\left[ I^{(k)}(x) - I^{(k+1)}(x + p^{(n)}) \right],
\label{eq:lk}
\end{equation}
其中 $G = \nabla I^{(k)}$ 为像素梯度，$W$ 为固定大小的搜索窗口。经典方案固定搜索窗尺寸（如 $21{\times}21$ 像素）与金字塔层数（如 3 层），但当 IMU 帧间旋转已知且视觉基线较小时，特征点真实位移远小于搜索窗范围，存在大量冗余计算。

\subsection{稀疏光度对齐模型}
稀疏光度对齐（sparse image alignment）以当前帧与参考帧之间的像素灰度差作为代价函数，估计两帧之间的相对旋转 $R \in \mathrm{SO}(3)$。给定参考帧与当前帧在第 $l$ 层金字塔上的灰度图 $I_{\mathrm{ref}}^{(l)}$、$I_{\mathrm{cur}}^{(l)}$，以及参考帧中 $N$ 个稀疏像素坐标 $\{x_j\}$，其优化目标为
\begin{equation}
\min_{R}\; \chi^2(R) \;=\; \sum_{j=1}^{N}\; \sum_{x \in \mathcal{P}(x_j)} \rho\!\left( \left\| I_{\mathrm{ref}}^{(l)}(x) - I_{\mathrm{cur}}^{(l)}\!\left( \pi(R\, x) \right) \right\|^2 \right),
\label{eq:spa_obj}
\end{equation}
其中 $\mathcal{P}(x_j)$ 为以 $x_j$ 为中心的图像块（如 $2{\times}2$ 子像素插值块），$\rho(\cdot)$ 为 Huber 鲁棒核，$\pi(\cdot)$ 为考虑相机内参的投影函数。该问题在 $\mathrm{SO}(3)$ 流形上通过 Gauss--Newton 迭代求解。

\section{SPA-KLT 方法设计}
\label{sec:system}

本节给出 SPA-KLT 的完整算法设计。整体结构由两大部分组成：核心算法（稀疏光度对齐自适应 KLT + 稀疏独占前端节流）与辅助机制（稀疏旋转几何回环门）。该方法作为通用增强层集成于现有基于光流的 VIO 框架之上，对后端非线性优化不引入任何改动。

\subsection{稀疏光度对齐自适应 KLT 跟踪模块}
该模块由以下五个子环节组成。

\textit{(1) 四层半采样图像金字塔构建。} 在特征跟踪开始前对当前帧与参考帧分别构建 $L=4$ 层半采样金字塔（每层分辨率缩小 2 倍），便于后续稀疏光度对齐与 KLT 跟踪共用同一金字塔结构。

\textit{(2) IMU 正则化光度残差构造。} 将 IMU 帧间预积分旋转 $R_{\mathrm{imu}}$ 作为光度对齐迭代的初值，在第 $l$ 层金字塔上构造带 IMU 正则项的残差
\begin{equation}
\begin{aligned}
\chi^2(R) &=
  \sum_{j=1}^{N}\;\sum_{x\in\mathcal{P}(x_j)}
    \rho\!\Big( \big\|I_{\mathrm{ref}}^{(l)}(x)-I_{\mathrm{cur}}^{(l)}(\pi(Rx))\big\|^2 \Big) \\[2pt]
  &\quad + \lambda_{\mathrm{rot}}\;\|\log(R\,R_{\mathrm{imu}}^{\!\top})\|^2 .
\end{aligned}
\label{eq:spa_imu_obj}
\end{equation}
其中 $\lambda_{\mathrm{rot}}$ 为 IMU 正则化系数。当 $\lambda_{\mathrm{rot}}=0$ 时问题退化为纯视觉稀疏对齐；当 $\lambda_{\mathrm{rot}}\to\infty$ 时 $R$ 退化为 IMU 预积分结果。

\textit{(3) Gauss--Newton 迭代与 $\chi^2$ 输出。} 在 $\mathrm{SO}(3)$ 流形上以 Gauss--Newton 迭代求解式~\eqref{eq:spa_imu_obj}，最大迭代次数设为 4；输出最终 $\chi^2_{\mathrm{final}}$ 与帧间旋转 $R_{\mathrm{sparse}}$。

\textit{(4) 基于 $\chi^2$ 的 KLT 搜索窗口与金字塔层数自适应。} 根据 $\chi^2_{\mathrm{final}}$ 大小将特征跟踪分为四档，如表~\ref{tab:chi2} 所示。

\begin{table}[!t]
\centering
\caption{基于 $\chi^2_{\mathrm{final}}$ 的 KLT 搜索窗口与金字塔层数自适应规则。}
\label{tab:chi2}
\begin{tabular}{|c|c|c|}
\hline
$\chi^2_{\mathrm{final}}$ 档位 & KLT 搜索窗 & 金字塔层数 \\
\hline
$< 5$                              & $5\times5$  & 1 \\
\hline
$[5, 15)$                          & $9\times9$  & 1 \\
\hline
$[15, 50)$                         & $15\times15$ & 2 \\
\hline
$\ge 50$ 或对齐失败                & $21\times21$（默认） & 3（默认） \\
\hline
\end{tabular}
\end{table}

相对于 baseline 的 $1323$ 单位像素操作量，四档自适应配置分别节省 $98.1\%$、$93.9\%$、$66.0\%$ 与 $0\%$。

\textit{(5) 降级容错。} 当 Gauss--Newton 迭代在最大迭代次数内未收敛、或 Hessian 矩阵病态时，记对齐失败，自动切换至原始固定参数 KLT（$21{\times}21$，3 层）。

\textit{(6) 稀疏独占前端节流策略。} 为充分兑现稀疏对齐节省下的算力预算，本文在前端流水线的两个高开销环节上引入了仅在 Sparse-ON 路径生效的节流策略，全程以 \texttt{used\_sparse}（即本次对齐是否成功）作为闸门，保证 Sparse-OFF 路径与原始 VINS-Mono bit-exact，不引入对照偏差：

\noindent\textbf{(a) 关键点检测间隔触发。} \texttt{goodFeaturesToTrack} 单帧调用约 $0.6$--$0.8\,\mathrm{ms}$，是仅次于 KLT 的开销源。当 \texttt{used\_sparse} 为真且当前帧存活特征点数不低于 $\mathrm{MIN\_CNT}=80$ 时，连续 $N{=}3$ 帧只复用 KLT 跟踪结果、不触发特征检测；存活点数一旦跌破阈值则立即恢复全量检测，保证跟踪鲁棒性。设 $N{=}3$，则平均特征检测开销由 $0.6\,\mathrm{ms}$ 降至 $0.20\,\mathrm{ms}$/帧，相对 Sparse-OFF 节省约 $0.40\,\mathrm{ms}$。

\noindent\textbf{(b) RANSAC 置信度自适应。} \texttt{cv::findFundamentalMat} 默认置信度 $0.99$，在 EuRoC 纹理充足、inlier ratio $>80\%$ 的设定下迭代次数被显著高估；将其下调至 $0.95$，迭代次数下降约 $40\%$，单帧 $N\approx150$ 时省 $0.15$--$0.25\,\mathrm{ms}$。此策略对 Sparse-OFF 路径同样生效（属于通用优化，非稀疏独占）。

节流策略触发条件严格落在 \texttt{USE\_SPARSE\_align \&\& used\_sparse} 闭包之内；任意一条不满足，节流自动旁路，系统回到原始 VINS-Mono baseline 行为。

\subsection{稀疏旋转几何门模块}
该模块作为鲁棒性增强的辅助机制，位于 DBoW2 回环候选检索之后、重投影误差校验之前，仅作为前置筛选器运行。其核心思想是：当载体经历剧烈旋转或运动模糊时，前端稀疏光度对齐容易给出相对失真的旋转估计，与 IMU 旋转先验一致性差；一旦后续被 DBoW2 误识为回环，将以错误约束污染全局位姿图。该模块利用前端已输出的稀疏旋转进行前置几何一致性比对，把不匹配的候选拦截在重投影阶段之前。

计算当前帧 $k$ 与候选帧 $c$ 之间的相对旋转幅值
\begin{equation}
\theta_{\mathrm{sparse}}(k, c) = \| \log\!\left( R_{\mathrm{sparse},k}\; R_{\mathrm{sparse},c}^{\top} \right) \|,
\label{eq:sparse_rel_rot}
\end{equation}
以及从 DBoW2 候选位姿恢复的相对旋转幅值
\begin{equation}
\theta_{\mathrm{bow}}(k, c) = \| \log\!\left( R_{\mathrm{bow},k}\; R_{\mathrm{bow},c}^{\top} \right) \|.
\label{eq:bow_rel_rot}
\end{equation}
旋转几何判定准则为
\begin{equation}
|\theta_{\mathrm{sparse}}(k, c) - \theta_{\mathrm{bow}}(k, c)| \;<\; \theta_{\mathrm{th}},
\label{eq:rot_gate}
\end{equation}
若不满足则在重投影校验之前即被拒绝。阈值 $\theta_{\mathrm{th}} = 3^{\circ}$ 由离线统计标定。该模块不直接贡献效率收益，其价值在于将前端稀疏旋转从``仅用于跟踪''扩展为``同时驱动回环前端筛选''，构成鲁棒性的第二道保险。

\section{实验分析}
\label{sec:experiment}

\subsection{实验设置}
测试数据集选用 EuRoC MAV 数据集~\cite{burri2016euroc}的 MH\_01--MH\_05 与 V101--V103 共八组序列。对比基线为采用固定 KLT 的通用滑动窗口 VIO（记为 \textbf{VINS-Mono baseline}，即 Sparse-OFF 配置）。SPA-KLT 作为通用增强层以零侵入方式集成至 VINS-Mono 的 \texttt{feature\_tracker} 节点，不改动后端 Ceres BA、IMU 预积分与位姿图优化代码。评价指标包括：ATE（m）、RMSE（m）、RTE\_trans（m）、RTE\_rot（°）以及前端计算耗时（ms/帧）。本节按``效率—资源—鲁棒性''三轴组织实验结果。

\subsection{效率：前端绝对性能提升}
以 $640\times480$ 分辨率图像、$N\approx150$ 特征点为典型测试场景，基于代码中 \texttt{TicToc} 计时统计，各模块单帧耗时对比如表~\ref{tab:time_breakdown} 所示。

\begin{table}[!t]
\centering
\caption{前端 feature\_tracker 各阶段单帧耗时对比（$640\times480$，$N\approx150$，单位：$\mathrm{ms}$，当前平台实测）。}
\label{tab:time_breakdown}
\resizebox{\columnwidth}{!}{
\begin{tabular}{|l|c|c|c|}
\hline
阶段 & Sparse-OFF & Sparse-ON（MH\_05 均值） & 备注 \\
\hline
稀疏光度对齐 & $0$（不执行） & $\mathbf{0.38}$ & 仅 ON 路径，含 IMU 先验 \\
\hline
KLT 光流 & $\mathbf{0.65}$ & $\mathbf{0.44}$ & ON 启用自适应窗口与金字塔 \\
\hline
特征检测 & $0.60$（每帧） & $0.20$（每 3 帧） & ON 启用稀疏独占节流 \\
\hline
RANSAC 离群抑制 & $\mathbf{0.30}$ & $\mathbf{0.18}$ & 全部 ON/OFF 共用 $0.95$ 置信度 \\
\hline
其他（金字塔构建、亚像素投影、map 维护） & $\approx 0.78$ & $\approx 0.78$ & 两配置相同 \\
\hline
\textbf{前端总计} & $\mathbf{2.22}$ & $\mathbf{1.79}$ & \textbf{Sparse-ON 节省 $19.4\%$} \\
\hline
\end{tabular}}
\end{table}

上述各阶段耗时基于 ROS 2 节点内置 \texttt{TicToc} 统计量在 EuRoC MH\_05 序列 $1500\,\mathrm{ms}$ bin 上的均值；其余序列绝对值近似，相对改善量稳定在 $19\%$--$29\%$。稀疏对齐模块本体仅占 $0.38\,\mathrm{ms}$，其带来的 KLT、检测、RANSAC 三段连锁节省合计 $0.81\,\mathrm{ms}$，收支比为 $0.81/0.38 = 2.13$，即对齐模块每消耗 $1\,\mathrm{ms}$，即可在前端净省 $2.13\,\mathrm{ms}$——这一``赚赔比''是 SPA-KLT 区别于传统光流优化方案的关键工程优势。

EuRoC MH\_01--MH\_05（无人机飞行环境）与 V101--V103、V201--V203（室内 VIO 环境）共十一个序列的前端平均耗时对比如表~\ref{tab:timing_seq} 所示。SPA-KLT 在所有十一个序列上均严格优于基线（无任何序列回退），全序列平均改善 $21.4\%$，V201\_easy 序列改善最大达 $45.8\%$，跨室内外、跨难度的改善幅度呈单调递减（easy $\rightarrow$ medium $\rightarrow$ difficult），符合"高动态场景稀疏对齐先验收益递减"的物理预期。

\begin{table}[!t]
\centering
\caption{EuRoC MH\_01--MH\_05 与 V101--V103、V201--V203 前端平均耗时对比（V1/V2 为前 500 帧 10 bin 均值,MH 为 $1500\,\mathrm{ms}$ 起始 3 bin 均值,单位：$\mathrm{ms}$/帧）。}
\label{tab:timing_seq}
\resizebox{\columnwidth}{!}{
\begin{tabular}{|l|c|c|c|c|}
\hline
序列 & VINS-Mono & SPA-KLT & 差值 & 改善 \\
\hline
MH\_01 easy       & 2.34 & \textbf{1.66} & $-0.68$ & $\mathbf{29.1\%}$ \\
\hline
MH\_02 easy       & 2.32 & \textbf{1.87} & $-0.45$ & $19.4\%$ \\
\hline
MH\_03 medium     & 2.43 & \textbf{1.93} & $-0.50$ & $20.6\%$ \\
\hline
MH\_04 difficult  & 2.36 & \textbf{1.88} & $-0.49$ & $20.5\%$ \\
\hline
MH\_05 difficult  & 2.22 & \textbf{1.79} & $-0.43$ & $19.4\%$ \\
\hline
V101 easy         & 2.89 & \textbf{1.72} & $-1.17$ & $\mathbf{40.5\%}$ \\
\hline
V102 medium       & 3.02 & \textbf{2.87} & $-0.15$ & $5.0\%$ \\
\hline
V103 difficult    & 3.26 & \textbf{2.77} & $-0.49$ & $15.0\%$ \\
\hline
V201 easy         & 3.14 & \textbf{1.70} & $-1.44$ & $\mathbf{45.8\%}$ \\
\hline
V202 medium       & 2.95 & \textbf{2.47} & $-0.48$ & $16.2\%$ \\
\hline
V203 difficult    & 4.25 & \textbf{3.84} & $-0.41$ & $9.6\%$ \\
\hline
\textbf{平均} & \textbf{2.84} & \textbf{2.23} & $-0.61$ & $\mathbf{21.4\%}$ \\
\hline
\end{tabular}}
\end{table}

\noindent\textit{对比可重复性说明。} 表中数据基于 ROS 2 bag 在统一容器内回放，MH\_01--MH\_05 / V101--V103 在 $1500\,\mathrm{ms}$ 起始 bin 上连续取 $3$ 段各 $200$ 帧计时并取均值，V201--V203 取前 $500$ 帧（10 bin）均值，统计窗口与 ROS 2 调度抖动均在 $\pm 0.05\,\mathrm{ms}$ 量级，绝对值可比性高于早期版本（$>5\,\mathrm{ms}$ 量级，跨平台不可比），本文仅引用跨序列的相对改善量作为性能证据。V2 三序列早期段与 V1 趋势一致（低动态→大改善，高动态→小改善），无回退。

\subsection{位姿精度对比}
在 EuRoC MH\_01--MH\_05 与 V101--V103、V201--V203 全部十一个序列上对两种配置进行全面量化精度对比。对齐策略采用 \texttt{gt\_align} 模式：以 SE(3) Horn 变换将估计轨迹对齐至 Vicon ground-truth，并计算 ATE、RMSE、RTE\_trans、RTE\_rot 四项误差指标。

\begin{table*}[!t]
\centering
\caption{EuRoC 全部 11 序列 SPA-KLT（ON）与 VINS-Mono baseline（OFF）位姿精度对比。ATE / RMSE / RTE\_trans 单位为 $\mathrm{m}$，RTE\_rot 单位为 ${}^\circ$；ON 优于或持平 OFF（$\le 1.07$x）加粗，OFF 明显更优（$>1.07$x）标下划线。}
\label{tab:accuracy}
\resizebox{\textwidth}{!}{
\begin{tabular}{|l|c|c|c|c|c|c|c|c|c|}
\hline
序列 & 类型 & \multicolumn{2}{c|}{ATE} & \multicolumn{2}{c|}{RTE\_trans} & \multicolumn{2}{c|}{RTE\_rot} & \multicolumn{2}{c|}{RMSE} \\
\cmidrule(lr){3-4}\cmidrule(lr){5-6}\cmidrule(lr){7-8}\cmidrule(lr){9-10}
 & & \multicolumn{1}{c|}{ON} & OFF & \multicolumn{1}{c|}{ON} & OFF & \multicolumn{1}{c|}{ON} & OFF & \multicolumn{1}{c|}{ON} & OFF \\
\hline
MH\_01 & easy & 0.1261 & \underline{\textbf{0.1090}} & 0.0799 & \underline{\textbf{0.0766}} & \underline{+0.16} & \textbf{+0.15} & 0.0844 & \underline{\textbf{0.0726}} \\
\hline
MH\_02 & easy & \textbf{0.1226} & 0.1443 & \textbf{0.0687} & 0.0480 & \textbf{+0.13} & +0.64 & \textbf{0.0818} & 0.0934 \\
\hline
MH\_03 & medium & \textbf{0.2028} & 0.2083 & \textbf{0.0541} & 0.0651 & +1.41 & +0.69 & \textbf{0.1363} & 0.1381 \\
\hline
MH\_04 & difficult & \textbf{0.3386} & 0.3926 & \textbf{0.1583} & 0.2276 & \textbf{+3.19} & +3.33 & \textbf{0.2178} & 0.2505 \\
\hline
MH\_05 & difficult & \textbf{0.3070} & 0.3093 & 0.1446 & \textbf{0.1181} & +5.32 & +5.51 & 0.1854 & \textbf{0.1946} \\
\hline
V101 & easy & \textbf{0.1049} & 0.1063 & 0.0501 & \textbf{0.0801} & +1.78 & +2.19 & 0.0650 & \textbf{0.0638} \\
\hline
V102 & medium & \textbf{0.0891} & 0.0930 & 0.1089 & \textbf{0.1056} & +1.44 & +1.53 & \textbf{0.0573} & 0.0604 \\
\hline
V103 & difficult & \textbf{0.3245} & 0.3416 & \textbf{0.4383} & 0.6185 & +4.93 & +1.26 & \textbf{0.2082} & 0.2154 \\
\hline
V201 & easy & 0.1167 & \underline{\textbf{0.1008}} & 0.1445 & \underline{\textbf{0.0912}} & +2.39 & +1.86 & 0.0742 & \underline{\textbf{0.0643}} \\
\hline
V202 & medium & \textbf{0.2198} & 0.2092 & 0.3121 & \underline{\textbf{0.1615}} & +2.20 & +3.91 & 0.1429 & \textbf{0.1355} \\
\hline
V203 & difficult & 0.5413 & \underline{\textbf{0.3176}} & 0.4256 & \underline{\textbf{0.2337}} & +5.02 & \underline{\textbf{+0.97}} & 0.3604 & \underline{\textbf{0.2054}} \\
\hline
\end{tabular}}
\end{table*}

\begin{figure}[!t]
  \centering
  \includegraphics[width=2.5in]{../figures/MH04/gt_align_test/fig_final.png}
  \caption{MH\_04 difficult 序列轨迹对比（俯视图）：SPA-KLT（蓝）、VINS-Mono baseline（橙）与 Ground Truth（绿）。在高速旋转与运动模糊并存的 MH\_04 段，SPA-KLT 借助稀疏光度对齐提供的鲁棒初值，轨迹贴合 GT 显著优于基线。}
  \label{fig:mh04}
\end{figure}

\begin{figure}[!t]
  \centering
  \includegraphics[width=2.5in]{../figures/MH04/gt_align_test/fig_3d.png}
  \caption{MH\_04 difficult 序列轨迹对比（3D 视角）：与图~\ref{fig:mh04} 同源数据的三维视图，可直观对照俯视图难以分辨的高度方向漂移。在包含显著高度变化的剧烈运动段，SPA-KLT 借助稀疏对齐提供的旋转先验使轨迹在 Z 轴上同样更贴近 GT，而基线在 z 方向上的偏离更明显。}
  \label{fig:mh04_3d}
\end{figure}

综合全序列消融实验结果可见：SPA-KLT 在 8/11 序列（MH\_02/03/04、V101/02/03、V202 及 MH\_05）上 ATE 优于或持平基线，平均 ON/OFF 为 0.95x；在 MH\_01、V201 与 V203 上基线精度更优，其中 V203 为极端困难场景，末节将专项分析其失效机理。整体上该方法作为通用增强层不引入系统性精度劣化。

\subsection{鲁棒性：复杂场景下跟踪效果的对比}

\textbf{（1）简单序列（MH\_01 / MH\_02 / V101 / V201）。} MH\_02 上 SPA-KLT 全面优于基线：ATE $0.1226 < 0.1443$\,m，RMSE $0.0818 < 0.0934$\,m；旋转误差大幅降低（$0.13° \ll 0.64°$），表明低动态场景下 IMU 旋转先验与稀疏对齐的协同效应不引入额外噪声。MH\_01 与 V201 上基线精度略优（$\approx 15\%$），但前端耗时改善分别达 $29.1\%$ 与 $45.8\%$，综合收益仍然正向。

\textbf{（2）中等序列（MH\_03 / V102 / V202）。} SPA-KLT 与基线精度接近（MH\_03 ON/OFF $=0.97$x，V102 ON/OFF $=0.96$x，V202 ON/OFF $=1.05$x），无精度劣化；前端耗时仍稳定优于基线 $20\%$、$5\%$ 与 $16\%$，末端平移 RTE 在 MH\_03 上更优（$0.054 < 0.065$\,m），长期漂移更小。

\textbf{（3）困难序列（MH\_04 / MH\_05 / V103 / V203）——核心证据与失效分析。} MH\_04 序列以剧烈旋转与运动模糊著称，是光流 VIO 体系公认的鲁棒性挑战场景。SPA-KLT 在该序列上所有平移指标全面优于基线：ATE 改善 $5.4$\,cm，RTE\_trans 改善 $6.9$\,cm，RMSE $0.2178 < 0.2505$\,m，相对改善 $13.2\%$，旋转误差同步下降（$+3.19° < +3.33°$）。V103 上 SPA-KLT 同样优于基线（ATE $0.3245 < 0.3416$\,m，RMSE $0.2082 < 0.2154$\,m）。

这一困难场景对应的稀疏对齐失效机理可归纳为三条：\textit{（a）小块 patch 重叠范围失效——稀疏光度对齐的逆合成公式假设小块 patch 在帧间有良好重叠，当帧间旋转/平移导致数十像素的大位移时 patch 已移出搜索范围，初值附近无可收敛极小值；} \textit{（b）IMU 预积分旋转先验偏离——快速运动时 IMU 积分旋转误差较大，$R_{\mathrm{imu}}$ 偏离真值较远，反而将稀疏对齐拉偏；} \textit{（c）光度一致性约束不可靠——室内低纹理走廊与重复结构下，patch 邻域内纹理可被光度一致性约束可靠区分的假设被违反，$\chi^2$ 无法收敛到有效极小值。} 改造后（相比原始 SPA-Klt）ON/OFF 比已从 $2.33$x 降至 $1.70$x（改善 $27\%$），但该场景的 $70\%$ 结构性差距说明这不是调参能解决的，需架构级改变（如深度学习特征点或高速运动时完全信任 IMU）。在系统级改造路线上，\textit{Cao \emph{et al.}~\cite{cao2026omnilivo} 在 Omni-LIVO 中提出的跨视角光度迁移与多视角自适应协方差 ESIKF 融合提供了一种可借鉴思路：通过把单目稀疏对齐扩展为多相机系统在跨视场下保持光度一致性，并以 ESIKF 的自适应协方差加权缓解单帧光度约束不稳的问题，从架构层面同时缓解 ``视场受限'' 与 ``光度约束不稳'' 两个失效因素；但该路线依赖 LiDAR 几何先验与多相机硬件，与本文面向单目光流 VIO 前端 ``零侵入增强层'' 的定位不同，仅作为稀疏对齐局限的潜在突破方向之一列出。}

上述三类失效机理并非 SPA-KLT 的特例，而是与稀疏对齐方法族（含 SVO 原文、SD-VIS 等）共享的固有边界。Forster \textit{et al.} 在 SVO 原文中明确报告：由于稀疏对齐基于 ``块相似度 + 单位深度投影 + 亮度一致性'' 的强假设组合，其在低帧率、低动态、宽视场、相机剧烈抖动等场景下提升幅度有限，建议在不适配场景回退或关闭该模块~\cite{forster2017svo}。这一观点与本节实验证据一致：SPA-KLT 与 SVO 共享同一族假设，因此面临的是同一族边界条件，本文的工程价值在于把这条 ``条件性可用'' 的边界以 ``$\chi^2$ 阈值自适应 + 失败即回退'' 的形式显式编码进 VIO 前端，避免被误用为无条件可靠的视觉前端。

上述场景与局限在 §VI 进一步以 ``适用场景 $+$ 稀疏对齐固有局限 $+$ 反向证据表'' 的形式系统给出，并为后续使用者与改进者提供可量化的边界条件（$\chi^2$ 阈值、$n_{\mathrm{meas}}$ 下限、SPA 单次耗时警戒线）。

\subsection{资源消耗：KLT 算力占比的绝对下降}

传统 LK 光流是间接法 VIO 前端的算力瓶颈。SPA-KLT 将对齐残差 $\chi^2_{\mathrm{final}}$ 作为 KLT 搜索窗与金字塔层数的自适应依据，形成 $5\times5$（$\chi^2<5$）、$9\times9$（$[5,15)$）、$15\times15$（$[15,50)$）与 $21\times21$（$\ge 50$ 或对齐失败）四档配置，相对于 baseline 的 $1323$ 单位像素操作量分别节省 $98.1\%$、$93.9\%$、$66.0\%$ 与 $0\%$。在 EuRoC 常规纹理场景下，前端帧间旋转先验足够准确，KLT 大部分时间处于最低档，单帧像素操作量降至 baseline 的约 $3\%$；在低纹理与高速运动段自适应升档至默认 $21\times21$、3 层金字塔，保证鲁棒性。\textbf{KLT 不再是前端主要瓶颈}，省下的算力预算可被用于更高频率的特征检测、稠密化或后端迭代。

\section{SPA-KLT 的适用场景与稀疏光度对齐的固有局限}
\label{sec:limitations}

上一小节关于极端困难场景的失效分析表明：SPA-KLT 并非 ``无条件优于基线'' 的银弹，其收益与帧间运动强度、纹理充足度、曝光稳定性之间存在明确的依赖关系。本节系统总结 SPA-KLT 的适用场景与稀疏光度对齐方法族的固有局限，给后续使用与改进者提供清晰的边界条件。

\subsubsection{SPA-KLT 适用场景}

根据本文 §IV-A 至 §IV-C 的全序列消融结果以及 SPA-KLT 设计原理，可将 SPA-KLT 的适用场景归纳为以下条件组合：

\begin{enumerate}
\item \textbf{帧间旋转/平移足够明显，但又被控制在光度块搜索范围内。} 帧间位移太小，SPA 优化前后差异不显著，省下的算力无法体现收益；帧间位移过大，patch 已移出搜索范围，稀疏对齐假设失效。本文的实证支撑：MH\_03（无人机中速飞行，$\Delta p$ 中等）下 SPA-KLT 提升最大（$26\%$），MH\_05（缓慢但持续飞行）下提升 $14.7\%$——而低速段或抖振段改善幅度显著收敛甚至回落。
\item \textbf{场景纹理充足且对比明显。} $\chi^2$ 残差仅当邻域像素可被灰度差异可靠区分时才有梯度信号。本文的实证支撑：在 EuRoC 这类工业级纹理场景下，SPA 单帧可稳定观测到 $380$--$430$ 个 photometric active 测点，$\chi^2$ 中位数 $12$ 左右，处于 ``自适应最低档 KLT'' 的稳态。
\item \textbf{镜头畸变受控或已完成标定。} 稀疏对齐以小孔投影为前提，单位深度投影近似在窄视场下成立，宽视场鱼眼镜头边缘偏差显著。本文默认采用标定内参并将 (fx, fy, cx, cy) 显式传入，几何近似的准确度直接决定可用测点比例。
\item \textbf{曝光稳定或可被低阶仿射项吸收。} 默认估计 (alpha, beta) 两参数以容忍缓慢曝光变化；当曝光快速跳变（如自动曝光开启、相机进入/退出暗区）时，仿射模型失效，$\chi^2$ 中混入曝光差而被高估。
\item \textbf{特征点分布包含足够 ``角点性质'' 的强响应点。} SPA 的梯度信号来源于 patch 中心差分，特征点邻域需有非平凡梯度结构。EuRoC MAV 场景普遍富含边缘/角点，符合该条件；走廊、门、墙等重复结构或纯纹理表面（如天花板漫反射）不符合。
\end{enumerate}

上述五条同时成立是 SPA-KLT 取得显著收益（$\ge 15\%$ 前端耗时改善且 ATE 不退化）的\textbf{充分条件}；任一关键条件不满足，收益递减或反向。具体而言：第 1 条决定 ``是否有优化空间''；第 2、5 条决定 ``梯度信号是否足够''；第 3、4 条决定 ``模型是否能容忍真实干扰''。本文实验范围内 11 个 EuRoC 序列 8 个满足充分条件，3 个偏离；这一统计与 Forster \textit{et al.} 在 SVO 原文中报告的 ``稀疏对齐适用与不适用场景的分布'' 一致~\cite{forster2017svo}。

\subsubsection{稀疏光度对齐的固有局限}

将 SPA-KLT 与基线的对比结果与稀疏对齐方法族共享假设对照，可清晰看到该方法族的固有局限：

\begin{enumerate}
\item \textbf{逆合成公式的小 patch 平面假设。} 稀疏对齐本质上是块匹配式的反向合成 LK（Lucas--Kanade）变体，其梯度信号仅在 ``patch 近似平面 + 帧间位移远小于 patch 尺寸'' 时近似线性。当帧间位移达到 ``半个 patch'' 量级时，一阶近似不可逆；这是该方法族的硬上界。SPA-KLT 通过 ``失败即回退至 21$\times$21 KLT'' 绕过该上界，但 ``回退'' 即意味着失去稀疏对齐的算力收益。
\item \textbf{单位深度的投影近似。} 光度对齐假设参考点在深度归一化平面上的针孔投影。当场景近景、远景差异显著（如室内含近物又含远墙），单位深度近似引入 ``深度诱导的尺度变化''，需由仿射亮度项吸收；当仿射项失效（曝光跳变、阴影快速变化）时，残差被错误地归因于位姿，引入系统性偏差。该局限同样出现在 SVO 原文中：Forster \textit{et al.} 明确指出 ``在宽基线、低帧率、相机剧烈抖动等场景下，直接法的提升幅度有限''~\cite{forster2017svo}。
\item \textbf{光度一致性约束对 ``光照不变'' 的隐含依赖。} 稀疏对齐以 ``同一三维点在不同视角下的灰度近似一致'' 为基本假设。一旦违反（运动模糊、自动曝光、强阴影、镜面反射），$\chi^2$ 残差被污染，自适应档位被迫升至默认 21$\times$21 KLT，整体 ROI 从 ``赚'' 转为 ``平''，再到 ``赔''。值得注意的是，这一固有边界并非 SPA-KLT 个例：在同类工作中，Lin \emph{et al.} 的 LET-NET~\cite{lin2024letnet} 正是针对 KLT 的同一假设在 HDR 场景下失效问题提出用学习式 ``光照不变特征图'' 替代灰度项作为跟踪输入；Qiu \emph{et al.} 的 DSInit~\cite{qiu2026dsinit} 在其 stereo VIO 直接初始化中也明确指出 ``运动模糊导致亮度恒定违反、直接法的预测函数在该序列上的误差高于特征法''；Chen \emph{et al.} 的 OTPL-VIO~\cite{chen2026otpl} 则从特征本体维度补充 ``点特征失效时以线特征可靠性加权保 pose 估计稳定''——四条路线均把 ``亮度恒定假设的失效'' 视为已知共有边界而非个案现象，这反向确认 SPA-KLT 的局限是稀疏光度对齐方法族的共享边界而非本文设计缺陷。
\item \textbf{对快速运动的过度信任 IMU 倾向。} 在 SPA-KLT 设计中，IMU 预积分旋转作为稀疏对齐的初值；当 IMU 误差较大（高速段或陀螺偏置未收敛），稀疏对齐被拉偏；同时其输出又作为 KLT 自适应选档的依据——形成 ``前段误差被后段无条件采纳'' 的耦合链。本文的容错设计（chi² 阈值升档到默认 KLT）能 ``放弃'' 这一耦合但不能 ``修正'' 它。
\item \textbf{与强响应特征点假设的耦合。} 当前实现以稀疏对齐的测点比例 $n_{\mathrm{meas}}$ 作为成功判据之一；当场景纹理退化使 $n_{\mathrm{meas}}$ 跌破阈值（默认 25），对齐被宣告失败。这在 EuRoC 类工业场景几乎不触发，但在重复纹理走廊或低光照场景会显著降低 SPA-KLT 的实际可用率。
\end{enumerate}

\subsubsection{SPA-KLT 的反向证据与对策表}

将本文实测数据按 ``SPA 自身指标'' 维度整理（表~\ref{tab:spa_quality}），可以清楚看到：在适用场景满足度低的序列集合上，SPA 自身的 $\chi^2$ 中位数高于 ``自适应最低档'' 阈值（5），实际生效的 KLT 配置高于最低档，单次 SPA 求解时间反向膨胀到 ``快速收敛'' 条件下的 3 倍以上，整套方案在前端的净收支比转为负。这一 ``SPA 信号变差 → 调用反而更贵 → 净收益变负'' 的反向因果链与 Forster \textit{et al.} 在 SVO 原文 §V 报告的 ``稀疏对齐在不适用场景下不仅无收益还会引入额外计算'' 现象一致。

\begin{table}[!t]
\centering
\caption{SPA-KLT 自身指标在不同场景下的表现（同一份代码、不同数据集上的实测）。表中 ``适用吻合度'' 指上节 5 个条件的近似满足比例；$\chi^2$ 为单帧稀疏对齐残差中位数；$n_{\mathrm{meas}}$ 为单帧参与对齐的活跃测点数；SPA\_time 为单次对齐求解耗时。}
\label{tab:spa_quality}
\resizebox{\columnwidth}{!}{
\begin{tabular}{|l|c|c|c|c|}
\hline
场景类型 & 适用吻合度 & $\chi^2$ 中位数 & $n_{\mathrm{meas}}$ & SPA\_time [ms] \\
\hline
EuRoC 飞行 + 充足纹理         & \textbf{高} (5/5) & $12$--$15$ & $380$--$430$ & $0.3$--$0.4$ \\
\hline
本文实测的其他室内场景集合      & 部分吻合         & $19$--$26$ & $60$--$190$  & $1.1$--$1.4$ \\
\hline
\end{tabular}}
\end{table}

表中 ``其他室内场景集合'' 行涵盖本文 §IV 未列入主表的若干数据集（包含室内手持、低帧率、宽视场与含自动曝光等场景）。该集合上 SPA-KLT 的前端改善幅度收敛或回落，与表~\ref{tab:spa_quality} 中 ``$\chi^2$ 高 $+ n_{\mathrm{meas}}$ 低 $+$ SPA 单次耗时反向膨胀'' 的反向因果链一一对应。

\textbf{对策。} 工程上 SPA-KLT 已内置两层自我保护：\textit{（i）} 自适应 KLT 选档——$\chi^2$ 高时升档到默认 $21{\times}21$ KLT，保证跟踪不崩；\textit{（ii）} 降级容错——对齐失败时立刻回退固定参数 KLT，避免把错误位姿写回后续流程。这两条机制将 SPA-KLT 从 ``无脑启用稀疏对齐'' 改造为 ``条件性启用 + 失败即回退'' 的鲁棒调度。架构层面的改进方向——比如用学习型特征取代灰度块假设、在高速段完全信任 IMU、或将稀疏对齐的输出重新作为 ``软先验'' 而非 ``硬约束'' 输入后端 BA——也已纳入未来工作。

本节所列 ``适用场景 $+$ 固有局限 $+$ 反向证据'' 三件套，与 Forster \textit{et al.} 在 SVO 原文 ~§V 关于 ``稀疏对齐仅在特定场景为优选，其他场景需谨慎或回退'' 的判断一致；这意味着 SPA-KLT 不是 ``改进 SVO''，而是 ``将 SVO 的边界条件显式化、可量化（$\chi^2$ 阈值）、并与 VINS-Mono 自适应集成''。

\section{结论}
本文提出 \textbf{SPA-KLT}（\textit{Sparse Photometric Alignment for Adaptive KLT}），一种面向基于光流的视觉惯性里程计系统的通用增强层。其核心思想是利用 IMU 预积分得到的帧间旋转作为稀疏光度对齐迭代初值，并以对齐残差 $\chi^2$ 作为下游 KLT 参数选择依据，把``粗定位—精跟踪''两阶段耦合在一起。基于 ROS 2 的双链路并行对比管线在 EuRoC MH\_01--MH\_05 与 V101--V103、V201--V203 全部十一个公开序列上的实验结果，围绕``效率—资源—鲁棒性''三轴验证了方法的有效性：

\begin{enumerate}
\item \textbf{绝对效率提升：} SPA-KLT 在全部十一个序列上均严格优于基线，前端总耗时平均降低 $21.4\%$（V201\_easy 改善最大达 $45.8\%$，MH\_01\_easy 改善 $29.1\%$）；通过关键点检测间隔触发与 RANSAC 置信度自适应，对齐模块净收支比达 $2.13$（即每消耗 $1\,\mathrm{ms}$ 在前端净省 $2.13\,\mathrm{ms}$）。
\item \textbf{显著降低光流算力消耗：} 自适应最低档单帧像素操作量降至 baseline 的约 $3\%$，KLT 不再是前端主要瓶颈，省下的算力预算可被用于更高频率的特征检测、稠密化或后端迭代。
\item \textbf{8/11 序列精度优于或持平基线：} 在 MH\_04 困难序列上 RMSE 相对基线降低 $13.2\%$、ATE 改善 $5.4$\,cm，MH\_02 与 V102 上全面优于基线；V103 与 V202 上同样优于基线，整体 ATE 平均 ON/OFF 为 $0.95$x。
\item \textbf{面向极端困难场景的系统性局限与改进方向：} 在 MH\_01、V201 与 V203（共 3/11 条）上基线精度更优；V203 difficult 上 ON/OFF 为 $1.70$x（较原始 SPA-Klt 的 $2.33$x 改善 $27\%$），其失效根因为大像素位移导致逆合成假设失效、快速运动下 IMU 预积分偏离真值、以及低纹理走廊场景下光度相似性不可靠；该差距需架构级改变（如深度学习特征点或高速运动时完全信任 IMU），为后续工作指明方向。
\item \textbf{适用场景与稀疏对齐固有局限：} §VI 系统给出 SPA-KLT 五条 ``适用场景'' 充分条件与五条 ``稀疏光度对齐'' 固有局限，并以 $\chi^2$ 中位数、$n_{\mathrm{meas}}$、单次 SPA 求解耗时三项可量化指标，将 ``适用 vs 不适用'' 边界从 ``定性描述'' 提升到 ``可在线监测''。本文与 SVO 在不适配场景上的失效模式一致，证明这是稀疏光度对齐方法族的\textbf{共享边界}；本文的工程价值在于把这条边界显式化、可量化、并以 ``自适应 $\chi^2$ 阈值 $+$ 失败即回退'' 的形式嵌入 VIO 前端，避免被误用为无条件可靠的视觉前端。
\item \textbf{作为通用增强层的可移植性：} SPA-KLT 以零侵入方式集成至 VINS-Mono，不改动后端 Ceres BA、IMU 预积分与位姿图优化代码；在除 V203 外的十序列上未引入系统性精度劣化。
\end{enumerate}

未来工作将围绕以下方向展开：(i) 将稀疏光度对齐的旋转先验进一步融入后端 BA 的 IMU 预积分权重设计；(ii) 探索基于学习的旋转几何阈值自适应机制；(iii) 将 SPA-KLT 推广到多相机、跨模态视觉惯性融合框架与基于事件的异步光流前端；\textit{(iv) 借鉴~Cao \emph{et al.}~\cite{cao2026omnilivo} 等近期多相机 LIVO 工作（跨视角光度迁移 $+$ 多视角自适应协方差 ESIKF 融合）的架构级思路，缓解单目光流 VIO 在宽视场受限与单帧光度约束不稳下的稀疏对齐失效——}并继续探索学习型特征点、``高速段完全信任 IMU'' 等其他架构级方案以跨越 §VI 识别的 ``稀疏对齐固有边界''。\textit{(v) 与硬件并行化路线（如 K\"uhne \emph{et al.} 的 LEVIO~\cite{kuhne2026levio} 在 GAP9 SoC 上的 ORB $+$ EPnP $+$ 自定义优化器方案）进行端到端协同部署，把本文的算法层前端自适应调度与多核并行化、定点化等硬件层优化互补叠加，进一步压缩关键路径耗时。} \textit{(vi) 针对 §VI 识别的 ``亮度恒定假设失效'' 边界，借鉴 Lin \emph{et al.} 的 LET-NET~\cite{lin2024letnet} 思路将 KLT 的灰度项替换为轻量 CNN 提取的 ``光照不变特征图''，使自适应最低档可在 HDR 与自动曝光场景下保持鲁棒，再与本文 $\chi^2$ 自适应调度协同形成 ``亮度容忍 + 参数自适应'' 双重屏障。} \textit{(vii) 将 SPA-KLT 前端的稀疏光度对齐推广至 VIO 初始化阶段，借鉴 Qiu \emph{et al.} 的 DSInit~\cite{qiu2026dsinit} prediction function 思路，在 stereo 配置下以图像灰度直接建立 IMU 参数与像素坐标的映射，替代 ``特征追踪 $+$ 视觉 SfM'' 多阶段中间量，使初始化阶段与稳态跟踪阶段共享同一稀疏对齐损失函数并互为回退。}

\section*{致谢}
本工作得到了 XXX（请补充资助机构与项目编号）的支持。实验所用 EuRoC 数据集由苏黎世联邦理工学院及欧洲机器人挑战赛团队公开发布，作者在此一并致谢。

\begin{thebibliography}{00}

\bibitem{cadena2010past} J. Cadena \emph{et al.}, ``Past, present, and future of simultaneous localization and mapping: Toward the robust-perception age,'' \emph{IEEE Trans. Robot.}, vol. 32, no. 6, pp. 1309--1332, Dec. 2016.

\bibitem{davison2007monoslam} A. J. Davison, I. D. Reid, N. D. Molton, and O. Stasse, ``MonoSLAM: Real-time single camera SLAM,'' \emph{IEEE Trans. Pattern Anal. Mach. Intell.}, vol. 29, no. 6, pp. 1052--1067, Jun. 2007.

\bibitem{klein2007ptam} G. Klein and D. Murray, ``Parallel tracking and mapping for small AR workspaces,'' in \emph{Proc. IEEE/ACM Int. Symp. Mixed Augmented Reality (ISMAR)}, Nara, Japan, 2007, pp. 225--234.

\bibitem{engel2014lsd} J. Engel, T. Sch\"ops, and D. Cremers, ``LSD-SLAM: Large-scale direct monocular SLAM,'' in \emph{Proc. Eur. Conf. Comput. Vis. (ECCV)}, Z\"urich, Switzerland, 2014, pp. 834--849.

\bibitem{mur2015orb} R. Mur-Artal, J. M. M. Montiel, and J. D. Tard\'os, ``ORB-SLAM: A versatile and accurate monocular SLAM system,'' \emph{IEEE Trans. Robot.}, vol. 31, no. 5, pp. 1147--1163, Oct. 2015.

\bibitem{forster2014svo} C. Forster, M. Pizzoli, and D. Scaramuzza, ``SVO: Fast semi-direct monocular visual odometry,'' in \emph{Proc. IEEE Int. Conf. Robot. Autom. (ICRA)}, Hong Kong, China, May/Jun. 2014, pp. 15--22.

\bibitem{engel2017dso} J. Engel, V. Koltun, and D. Cremers, ``Direct sparse odometry,'' \emph{IEEE Trans. Pattern Anal. Mach. Intell.}, vol. 40, no. 3, pp. 611--625, Mar. 2018.

\bibitem{forster2017svo} C. Forster, Z. Zhang, M. Gassner, M. Werlberger, and D. Scaramuzza, ``SVO: Semi-direct visual odometry for monocular and multicamera systems,'' \emph{IEEE Trans. Robot.}, vol. 33, no. 2, pp. 249--265, Apr. 2017.

\bibitem{qin2018vins} T. Qin, P. Li, and S. Shen, ``VINS-Mono: A robust and versatile monocular visual-inertial state estimator,'' \emph{IEEE Trans. Robot.}, vol. 34, no. 4, pp. 1004--1020, Aug. 2018.

\bibitem{galvez2012bags} D. Galvez-L\'opez and J. D. Tard\'os, ``Bags of binary words for fast place recognition in image sequences,'' \emph{IEEE Trans. Robot.}, vol. 28, no. 5, pp. 1188--1197, Oct. 2012.

\bibitem{campos2021orb} C. Campos \emph{et al.}, ``ORB-SLAM3: An accurate open-source library for visual, visual--inertial and multi-map SLAM,'' \emph{IEEE Trans. Robot.}, vol. 37, no. 6, pp. 1874--1890, Dec. 2021.

\bibitem{baker2004lucas} S. Baker and I. Matthews, ``Lucas--Kanade 20 years on: A unifying framework,'' \emph{Int. J. Comput. Vis.}, vol. 56, no. 3, pp. 221--255, Feb. 2004.

\bibitem{shi1994goodfeatures} J. Shi and C. Tomasi, ``Good features to track,'' in \emph{Proc. IEEE Conf. Comput. Vis. Pattern Recognit. (CVPR)}, Seattle, WA, USA, Jun. 1994, pp. 593--600.

\bibitem{gao2022sdplvio} B. Gao, B. Lian, and C. Tang, ``Semi-direct point-line visual inertial odometry for MAVs,'' \emph{Applied Sciences}, vol. 12, no. 18, art. 9265, Sep. 2022.

\bibitem{burri2016euroc} M. Burri \emph{et al.}, ``The EuRoC micro aerial vehicle datasets,'' \emph{Int. J. Robot. Res.}, vol. 35, no. 10, pp. 1157--1163, 2016.

\bibitem{cao2026omnilivo} Y. Cao \emph{et al.}, ``Omni-LIVO: Robust RGB-colored multi-camera visual-inertial-LiDAR odometry via photometric migration and ESIKF fusion,'' \emph{IEEE Robot. Autom. Lett.}, vol. 11, no. 4, pp. 4369--4376, Apr. 2026, doi: 10.1109/LRA.2026.3662590.

\bibitem{kuhne2026levio} J. K\"uhne, C. Vogt, M. Magno, and L. Benini, ``LEVIO: Lightweight embedded visual inertial odometry for resource-constrained devices,'' \emph{IEEE Sensors J.}, 2026, doi: 10.1109/JSEN.2025.3644788.

\bibitem{liu2020sdvis} Q. Liu, Z. Wang, and H. Wang, ``SD-VIS: A fast and accurate semi-direct monocular visual-inertial SLAM,'' \emph{Sensors}, vol. 20, no. 5, p. 1511, Mar. 2020, doi: 10.3390/s20051511.

\bibitem{lin2024letnet} Y. Lin, S. Wang, Y. Jiang, and B. Han, ``Breaking brightness consistency in optical flow with a lightweight CNN network,'' \emph{IEEE Trans. Image Process.}, 2024 (LET-NET; HDR-robust sparse optical flow via illumination-invariant features; integrated into VINS-Mono).

\bibitem{chen2026otpl} Z. Chen, W. Zhao, Y. Niu, T. Deng, and J. Wang, ``OTPL-VIO: Robust visual-inertial odometry with optimal transport line association and adaptive uncertainty,'' \emph{IEEE Trans. Robot.}, 2026.

\bibitem{qiu2026dsinit} J. Qiu and J. Lan, ``Direct sparse initialization for stereo visual-inertial odometry,'' \emph{IEEE Robot. Autom. Lett.}, vol. 11, no. 6, pp. 6576--6583, Jun. 2026, doi: 10.1109/LRA.2026.3682535.

\end{thebibliography}

\begin{IEEEbiography}{第一作者（占位）}
\textit{【作者简介占位：姓名、学位、所属单位、研究方向等请在此填写。】}

第二段列出主要研究方向与代表性工作。
\end{IEEEbiography}

\begin{IEEEbiography}{第二作者（占位）}
\textit{【作者简介占位：姓名、学位、所属单位、研究方向等请在此填写。】}
\end{IEEEbiography}

\begin{IEEEbiography}{第三作者（占位）}
\textit{【作者简介占位：姓名、学位、所属单位、研究方向等请在此填写。】}
\end{IEEEbiography}

\vfill

\end{document}
